%!TEX root = main.tex

\section{Communication Upper Bounds}\label{sec:communication-upper}

In this section, we will see a communication upper bound for computing
the degeneracy of a graph exactly.
The setting is as follows:
Alice and Bob are respectively given \emph{disjoint} edge sets $E_A$ and $E_B$ on the set of
vertices $V$ (known to both Alice and Bob), and wish to compute the degeneracy of the graph
$G = (V, E_A \cup E_B)$.
The main theorem of this section is:

\begin{theorem}\label{thm:communication-upper}
  There exists a deterministic communication protocol that, given an $n$-vertex
  graph $G$, computes the degeneracy $\kappa(G)$ using $O(n \log^3 n)$ bits of
  communication.
  Moreover, the protocol outputs a $\kappa(G)$-ordering of the vertices
  of $G$, and a $\kappa(G)$-core in $G$.
\end{theorem}

We will focus on a decision version of the degeneracy problem, which we define
in the lemma below:

\begin{lemma}\label{lem:degen-decision}
  There exists a deterministic communication protocol that given an $n$-vertex
  graph $G$ and an integer $k$, decides whether the degeneracy
  $\kappa(G) \leq k$ using $O(n \log^2 n)$ bits of communication.
  Moreover, if the protocol accepts, it outputs a $k$-ordering of the vertices
  of $G$, and if it rejects, it outputs a $(k + 1)$-core in $G$.
\end{lemma}

For brevity, we will call the decision version of degeneracy defined in the
lemma above $\degen_{n, k}$.
Note that once we have a protocol for the decision problem, a protocol for the
search variant follows via binary search.
\begin{proof}[Proof of \Cref{thm:communication-upper}]
  Alice and Bob can binary search for the minimum $k$ such that $G$ has
  degeneracy at most $k$ (which is exactly $\kappa(G)$) using the protocol
  given by \Cref{lem:degen-decision}.
  Note that this incurs a multiplicative $\log n$ overhead in communication
  cost.
  Further, the protocol in the lemma also gives a $\kappa(G)$-ordering, since
  it accepts the input $(G, \kappa(G))$.
  Finally, to obtain a $\kappa(G)$-core, Alice and Bob can run the protocol
  in the lemma on the input $(G, \kappa(G) - 1)$.
\end{proof}

In the rest of this section, we will prove \Cref{lem:degen-decision}.
First, we show an easier (but more expensive) protocol that captures the main
idea.

\subsection{Warm-Up: An \texorpdfstring{$\Ot{n\sqrt{n}}$}{Õ(n sqrt n)} Communication Protocol}\label{sec:nrootn}
Recall the classical peeling algorithm, where one successively peels off the
vertex of minimum degree to find a degeneracy ordering of $G$.
We state a decision version of it below, which removes an arbitrary vertex of
degree $\leq k$ in each step.

\begin{Algorithm}\label{alg:peeling}
  The Peeling Algorithm. \medskip

  \textbf{Input:} A graph $G$, and an integer $k$.\\
  \textbf{Output:} \true{} if $G$ has degeneracy $\leq k$, \fail{} otherwise.

  \begin{enumerate}
    \item While $G$ is non-empty:
      \begin{enumerate}
        \item If there is a vertex $v \in G$ of degree $\leq k$,
          remove $v$ from $G$.
        \item Else output \fail.
      \end{enumerate}
    \item Output \true.
  \end{enumerate}
  %  \caption{The Peeling Algorithm for $\degen(G, k)$}\label{alg:peeling}
\end{Algorithm}

Note that in the \true{} case the $k$-ordering is simply the order the vertices
were peeled in, and in the \fail{} case, the $(k + 1)$-core is the set of vertices
remaining in $G$ when there are no vertices of degree $\leq k$.

We want to somehow port this algorithm to the communication setting.
The immediate problem is that Alice or Bob cannot compute a low degree
vertex without any communication after a vertex is removed.
The naive way to get around this is to send the updated degree of each
vertex after each deletion, but this can cost $\Omega(n^2)$ communication.
Instead, we observe that we only need to update the degrees of a remaining
vertex if it is not too large:

\begin{observation}\label{obs:sqrt-local}
  If $\deg(u) > k + \sqrt{n}$ at some step of the peeling algorithm,
  then $u$ cannot be deleted in the next $\sqrt{n}$ steps, since its degree
  falls by at most $1$ per step.
\end{observation}

Hence we have the following algorithm:
\begin{enumerate}
  \item Compute the set $S = \Set{u \in V \given \deg(u) > k + \sqrt{n}}$ of
    \emph{safe} vertices, and set it aside.
  \item Perform $\sqrt{n}$ iterations of the peeling algorithm, while updating
    the degree of a neighbor $u$ of a deleted vertex only if $u \notin S$ and $\deg(u) > k$.
  \item If the graph is non-empty, go to step 1.
\end{enumerate}

Let $N_A(u)$ and $\deg_A(u)$ denote the neighborhood and degree of $u$ within
$(V, E_A)$ (and similarly $N_B(u)$ and $\deg_B(u)$ for $E_B$).
We provide an implementation of the algorithm above in a communication protocol;
the sets $D$ and $L$ hold vertices that are ready to \emph{delete}, and
\emph{low} degree respectively.

\begin{Algorithm}\label{alg:nrootn}
  The $\Ot{n\sqrt{n}}$ communication protocol for $\degen_{n, k}$. \medskip

  \textbf{Input:} Alice and Bob get disjoint edge sets $E_A$ and $E_B$
    respectively, and both players get an integer $k$.\\
  \textbf{Output:} \true{} if $G = E_A \cup E_B$ has degeneracy $\leq k$, \fail{} otherwise.
  \begin{enumerate}
    \item Alice and Bob send $\deg_A(u)$ and $\deg_B(u)$ respectively for every
      vertex $u \in V(G)$, so they can then compute:
      \begin{itemize}
        \item $D \gets \Set*{ u \in G \given \deg(u) \leq k }$.
        \item $L \gets \Set*{ u \in G \given k + 1 \leq \deg(u) \leq k + \sqrt{n} }$.
        \item $S \gets V(G) \setminus (D \cup L)$.
      \end{itemize}
    \item For $\sqrt{n}$ iterations:
      \begin{enumerate}
        \item If $D$ is empty, output \fail.
        \item Let $v$ be any vertex in $D$. To remove $v$ from $G$:
        \item\label{step:Lnbrs} Alice and Bob send $N_A(v) \cap L$ and $N_B(v) \cap L$ respectively.
        \item\label{step:Lupdt} Hence they can update $\deg(u)$ for all $u \in L$.
        \item\label{step:Dmove} If $\deg(u) < k + 1$ for any $u \in L$,
          Alice and Bob move it to $D$.
      \end{enumerate}

    \item If $G$ is non-empty, go to Step 1.
    \item Output \true.
  \end{enumerate}
\end{Algorithm}
Note that the proofs (a $k$-ordering of the vertices in the \true{} case,
and a $(k + 1)$-core in the \fail{} case) fall out of the protocol exactly as before.
The correctness of this protocol follows from \Cref{obs:sqrt-local}; we omit
the proof here since it is subsumed by the next section.
To see that it only uses $\Ot{n \sqrt{n}}$ bits of communication, we observe
that:
\begin{itemize}
  \item Sending the degree of all vertices costs $O( n \log n )$ bits, and
    we do this at most $\sqrt{n}$ times.
  \item Once a vertex is in $D$, no communication is spent for updating its own degree.
  \item Each vertex in $L$ features at most $\sqrt{n}$ times in the neighborhood
    of a deleted vertex before it is moved to $D$, costing $O( \sqrt{n} \log n )$
    bits per vertex.
\end{itemize}
Hence the total communication cost of the protocol is $O( n\sqrt{n} \log n)$.

\subsection{An \texorpdfstring{$\Ot{n}$}{Õ(n)} Communication Protocol}

There are two changes from the previous protocol:
\begin{enumerate}
  \item 
    Instead of a single safe set of vertices, we have $O( \log n )$ safe sets
    corresponding to exponentially increasing gaps between the degree and $k$.
    More precisely, we define the set $S_i$ for $i \in \brac{\log n}$ such that
    whenever Alice and Bob compute $\deg(u)$ for some vertex $u$, they place it
    in $S_i$ if $k + 2^{i - 1} \leq \deg(u) < k + 2^i$.
  \item
    Instead of counting a global number of iterations until the safe vertices
    are updated, we track each vertex individually.
    Concretely, suppose $\deg(u) \geq k + \ell$.
    Then $u$ cannot be deleted until $\deg_A(u)$ or $\deg_B(u)$ falls by
    at least $\ell / 2$.
\end{enumerate}
The upshot of the two changes above is that for $u \in S_i$ Alice and Bob do not
have to update $\deg(u)$ until $\deg_A(u)$ or $\deg_B(u)$ falls by at least
$2^{i - 2}$.
We can now state the algorithm.
As before, the set $D$ holds all vertices which are ready for deletion.

\begin{Algorithm}\label{alg:n}
  The $\Ot{n}$ communication protocol for $\degen_{n, k}$. \medskip

  \textbf{Input:} Alice and Bob get disjoint edge sets $E_A$ and $E_B$
  respectively, and both players get an integer $k$.\\
  \textbf{Output:} \true{} if $G = E_A \cup E_B$ has degeneracy $\leq k$, \fail{} otherwise.
  \begin{enumerate}
    \item Alice and Bob send $\deg_A(u)$ and $\deg_B(u)$ respectively for every
      vertex $u \in V(G)$, so they can then compute:
      \begin{itemize}
        \item $D \gets \Set*{ u \in G \given \deg(u) \leq k }$.
        \item $S_i \gets \Set*{ u \in G \given k + 2^{i - 1} \leq \deg(u) < k + 2^i }$,
          for each $i \in \brac{\log n}$.
      \end{itemize}
    \item While {$G$ is non-empty:}
      \begin{enumerate}

        \item If $D$ is empty, output \fail.
        \item Let $v$ be any vertex in $D$. To remove $v$ from $G$:
        \item For each $i \in \brac{ \log n }$, and each $u \in S_i$,
          if $\deg_A(u)$ or $\deg_B(u)$ has fallen by $\geq 2^{i - 2}$ since it
          was last communicated, Alice sends $\deg_A(u)$ and Bob sends $\deg_B(u)$,
          hence allowing them to update $\deg(u)$.
        \item For all vertices updated in the previous step, Alice and Bob place
          them in the correct $S_i$ (or $D$).
      \end{enumerate}
    \item Output \true.
  \end{enumerate}
  %  \caption{The Peeling Algorithm for $\degen(G, k)$}\label{alg:peeling}
\end{Algorithm}

To wrap up, we prove that the algorithm above is correct, and uses $\Ot{n}$
space.

\begin{proof}[Proof of \Cref{lem:degen-decision}]
To prove that this protocol is correct, we argue that:
\begin{itemize}
  \item All vertices have degree $\leq k$ when they are removed, which
    implies that if the protocol outputs \true{}, $\kappa(G) \leq k$.
    This is immediate because we only move vertices to $D$ when we know their
    current degree exactly, and it is $\leq k$.
  \item Any vertex $u$ with degree $\leq k$ is placed in $D$, which implies
    that if $\kappa(G) \leq k$, then the protocol outputs \true{}.
    Suppose towards a contradiction that $u$ was placed in the set $S_i$ when
    its degree was updated for the final time.
    This means that $\deg(u)$ was at least $k + 2^{i - 1}$.
    But now for $\deg(u)$ to fall below $k + 1$, at least one of $\deg_A(u)$
    or $\deg_B(u)$ must fall by $\geq 2^{i - 2}$ (otherwise $\deg(u)$ falls by
    less than $2^{i - 1}$).
    Hence $\deg(u)$ was updated again following the final update, and we have
    a contradiction.
\end{itemize}

Next, we argue that this protocol uses only $\Ot{n}$ communication.
\begin{itemize}
  \item As before, the initial setup uses only $O(\log n)$ bits per vertex.
  \item While in $S_i$, a vertex can have its degree updated at most twice
    before it falls below $k + 2^{i - 1}$, and is hence dropped from
    $S_i$.
    Since there are only $\log n$ sets $S_i$, and a vertex can never gain
    degree, we only update its degree $O(\log n)$ times before it falls into $D$,
    spending $O( \log^2 n )$ bits overall.
\end{itemize}
Hence, in total, the protocol uses only $O(n \log^2 n)$ bits of communication.
\end{proof}
